\documentclass{article}

\PassOptionsToPackage{numbers, sort, compress}{natbib}
\usepackage{neurips_2026}
\nolinenumbers

\usepackage[T1]{fontenc}
\usepackage[utf8]{inputenc}
\usepackage{hyperref}
\usepackage{url}
\usepackage{booktabs}
\usepackage{nicefrac}
\usepackage{microtype}
\usepackage{xcolor}
\usepackage{caption}

%%%

\usepackage{subcaption}
\usepackage{graphicx}
\usepackage{multirow}
\usepackage{amsmath,amssymb,amsfonts}
\usepackage{amsthm}
\usepackage{mathrsfs}
\usepackage{enumitem}
\usepackage{algorithm}
\usepackage{algorithmicx}
\usepackage{algpseudocode}

\newtheorem{theorem}{Theorem} % continuous numbers
%%\newtheorem{theorem}{Theorem}[section] % sectionwise numbers
%% optional argument [theorem] produces theorem numbering sequence instead of independent numbers for Proposition
\newtheorem{proposition}[theorem]{Proposition}% 
\newtheorem{lemma}{Lemma}% 
%%\newtheorem{proposition}{Proposition} % to get separate numbers for theorem and proposition etc.

\newtheorem{example}{Example}
\newtheorem{remark}{Remark}

\newtheorem{definition}{Definition}
\newtheorem{assumption}{Assumption}

%%%

\newcommand{\R}{\mathbb{R}}
\newcommand{\cN}{\mathcal{N}}
\newcommand{\cL}{\mathcal{L}}
\newcommand{\cK}{\mathcal{K}}
\newcommand{\cD}{\mathcal{D}}
\newcommand{\cF}{\mathcal{F}}
\newcommand{\KL}{\mathrm{KL}}
\newcommand{\norm}[1]{\left\|#1\right\|_w}
\newcommand{\normsq}[1]{\left\|#1\right\|_w^{2}}
\newcommand{\ip}[2]{\langle #1,\, #2\rangle}
\newcommand{\hats}{\hat{s}}
\newcommand{\T}{\mathcal{T}}
\newcommand{\cG}{\mathcal{G}}


\title{TEMPO: Regime-Aware Operator Learning with Locally Adapted POD Bases}

\author{Anonymous Authors}





\begin{document}


\maketitle

\begin{abstract}
Operator learning methods typically approximate the solution map of a physical system with a single global model — an assumption that fails when solutions change character fundamentally across the parameter space.

We propose TEMPO: a framework that identifies latent dynamical regimes in snapshot data via Expectation-Maximization, builds a dedicated NeuralPOD basis for each, and learns to route new inputs to the right regime at inference time. On multi-regime benchmarks, TEMPO achieves substantially lower cross-regime errors than per-parameter specialists, with interpretable structure that reflects the underlying physics.

\end{abstract}

\section{Introduction}\label{sec:intro}

Operator learning approximates the solution map of a parametric PDE from data,
enabling instant prediction without re-solving the governing equations at each
new parameter value~\cite{Kovachki2024}.
The dominant paradigm — exemplified by DeepONet~\cite{Lu2019} and
POD-DeepONet~\cite{Lu2022} — trains a single global model over the entire
parameter space.
This works well when solutions form a simple manifold, but breaks down when
the governing parameter drives \emph{qualitatively distinct} dynamics:
viscosity separating laminar diffusion from near-discontinuous shocks, or source
magnitude producing solutions that differ by orders of magnitude in amplitude.
In these multi-regime settings, a single global basis must compromise between
incompatible geometries, concentrating error precisely where accurate prediction
matters most.

Several methods improve expressivity while retaining a single global
representation.
POD-PoU~\cite{Sharma2025} trains a partition-of-unity mixture of POD bases
with a shared branch, reducing error on sharp-gradient problems but without
discovering the underlying regime structure.
NeuralPOD~\cite{Mou2026} constructs nonlinear orthogonal basis functions via
greedy residual minimization, representing each mode as a neural network and
thus overcoming the linear subspace constraint.
Both, however, operate on the full dataset globally and provide no operator
for routing new inputs to the appropriate representation at inference.

Clustering snapshots and fitting local bases to each cluster is a classical
idea in reduced-order modeling~\cite{Amsallem2012, Hess2019}.
Daniel et al.~\cite{Daniel2020} extend this to deep learning by clustering
solution subspaces on the Grassmann manifold and training a dictionary
selector at inference.
These methods use hard cluster assignments, operate offline on fixed datasets,
and do not couple regime discovery with the operator learning objective —
leaving basis quality and assignment accuracy disconnected.

We propose \textbf{TEMPO} (\textit{Regime-Aware Operator Learning with
Locally Adapted POD Bases}), which closes this gap in a principled
probabilistic framework.
TEMPO models the trajectory ensemble as a Gaussian mixture and recovers regime
structure via Expectation-Maximization (EM): the E-step softly assigns each
trajectory to regimes based on reconstruction quality, while the M-step
rebuilds a dedicated NeuralPOD basis per regime — so as bases sharpen,
assignments improve, and vice versa.
After convergence, a GatingNet and per-regime BranchNets are trained jointly,
routing new inputs to the correct regime at inference.

\textbf{Contributions:}
\begin{itemize}[leftmargin=1.5em]
  \item A Gaussian mixture model with EM that discovers $M$ latent regimes
    and fits a dedicated NeuralPOD basis per regime, requiring no regime labels.
  \item NeuralPOD-DeepONet: the first deployable operator built on the
    NeuralPOD basis of~\cite{Mou2026}, pairing it with a branch network that
    predicts mode coefficients for unseen inputs.
  \item An adaptive basis update rule based on a distribution-shift criterion,
    which determines when each regime's basis warrants retraining and avoids
    redundant computation across EM iterations.
  \item On 1D Burgers' and 2D Darcy flow, a single TEMPO model trained jointly
    across all parameter values substantially outperforms single-regime
    baselines, with discovered regimes that reflect intrinsic solution geometry
    rather than parameter labels.
\end{itemize}


\section{Problem Statement}\label{sec:problem}

Let $\Omega \subset \R^d$ be a bounded spatial domain and $\T \subset \R_{\geq 0}$ a time interval.

We are given a dataset of $N$ solution trajectories:
\begin{equation}
  \cD = \bigl\{\,(u^{(i)}_0,\,\kappa^{(i)},\,s^{(i)})\,\bigr\}_{i=1}^{N},
  \qquad
  s^{(i)} = s\!\left(\cdot,\,\cdot\,;\,u^{(i)}_0,\kappa^{(i)}\right)
  \in L^2(\Omega\times\mathcal{T}),
  \label{eq:dataset}
\end{equation}
where $u_0^{(i)} \in \mathcal{U} = L^2(\Omega)$ is the initial condition and $\kappa^{(i)} \in \cK$ the physical parameter.
Each trajectory $s^{(i)}$ is discretized on $N_y = N_t N_x$ quadrature points
$\{y_j = (x_j, t_j),\, w_j\}_{j=1}^{N_y}$ covering the product grid
$\Omega \times \mathcal{T}$, with weighted norm $\|f\|_w^2 = \sum_j w_j f(y_j)^2$.

\begin{definition}[Solution operator]
\label{def:operator}
The \emph{solution operator}
 \[ \cG: \mathcal{U} \times \cK \to L^2(\Omega \times \T) \]
maps each pair $(u_0,\kappa)$ to the full space-time solution:
$\cG(u_0,\kappa) = s(\,\cdot\,,\,\cdot\,;\,u_0,\kappa)$.
\end{definition}


The operator $\cG$ is unknown and expensive to evaluate directly. The \emph{operator learning problem} is to construct $\hat{\cG} \approx \cG$ from $\cD$ alone, generalizing to unseen $(u_0,\kappa) \notin \cD$ without solving the PDE.


\paragraph{Multi-regime assumption.}

When solutions change character fundamentally across $\cK$, the optimal low-dimensional representation differs from regime to regime, and no single global basis captures the full solution diversity. We formalize this observation via a latent variable model over the trajectory ensemble.

\begin{definition}[Generative model]
\label{def:generative}
Each trajectory arises from one of $M$ latent regimes. The regime label $z_i \in \{1,\ldots,M\}$ is unobserved and drawn from a categorical prior,
\begin{equation}
  z_i \sim \mathrm{Categorical}(\pi_1,\ldots,\pi_M),
  \qquad \pi_m \geq 0,\quad \textstyle\sum_m \pi_m = 1.
  \label{eq:prior}
\end{equation}
Conditioned on $z_i = m$, the trajectory is distributed as
\begin{equation}
 ( s^{(i)} \mid z_i{=}m) \;\sim\;
  \cN\!\bigl(\hats_m^{(i)},\;\sigma^2 W^{-1}\bigr),
  \label{eq:likelihood}
\end{equation}
where $W = \mathrm{diag}(w_1,\ldots,w_{N_y})$ is the quadrature weight matrix, $\sigma^2 > 0$ controls the softness of regime assignments, and $\hats_m^{(i)}$ is the regime-$m$ approximation to $s^{(i)}$.
\end{definition}


The generative model leaves $\hats_m^{(i)}$ unspecified. We require only that it admits a \emph{low-rank trajectory decomposition}:
\begin{equation}
  \hats_m^{(i)}(y)
  \;=\; \phi_0^{(m)}(y)
  + \sum_{k=1}^{K_m}
    \lambda_k^{(m,i)}\,\phi_k^{(m)}(y),
  \qquad y = (x,t) \in \Omega\times\mathcal{T},
  \label{eq:separable}
\end{equation}
where $\phi_k^{(m)} : \Omega\times\mathcal{T}\to\mathbb{R}$ are spatiotemporal
mode functions and $\lambda_k^{(m,i)} \in \mathbb{R}$ a scalar amplitude
for trajectory~$i$ on mode~$k$.
This form encompasses POD~\citep{Lu2022} and learned neural bases~\citep{Mou2026};
in this work we adopt the latter.

\begin{definition}[Regime approximator]
\label{def:approx}
The parameters of regime $m$ are
$\Phi_m = \{\phi_k^{(m)}\}_{k=0}^{K_m}$
and scalar amplitudes
$\Lambda_m = \{\lambda_k^{(m,i)}\}_{k=1,\,i=1}^{K_m,\,N}$,
where $\phi_0^{(m)} : \Omega\times\mathcal{T}\to\mathbb{R}$ is the regime mean field,
$\{\phi_k^{(m)}\}_{k \geq 1}$ spatiotemporal mode functions,
$\lambda_k^{(m,i)}\in\mathbb{R}$ the scalar amplitude of trajectory~$i$ on mode~$k$,
and $K_m \leq K_{\max}$ the adaptive mode count.
\end{definition}

The learning problem decomposes into two phases: discovering the latent regime structure from $\cD$ \emph{offline}, then training a deployable operator for fast inference \emph{online}.


\paragraph{Phase 1: Offline regime discovery.}

Since the regime labels $z_i$ in~\eqref{eq:likelihood} are unobserved,
we recover the bases $\{\Phi_m\}$, mixture weights $\pi$, and noise
level $\sigma^2$ by maximizing the observed-data log-likelihood:
\begin{equation}
  (\Phi^*,\pi^*,\sigma^{2*})
  \;=\; \operatorname*{arg\,max}_{\Phi,\,\pi,\,\sigma^2}
  \;\sum_{i=1}^N \log \sum_{m=1}^M
    \frac{\pi_m\,\det(W)^{1/2}}{(2\pi\sigma^2)^{N_y/2}}
    \exp\!\!\left(-\frac{\normsq{s^{(i)}-\hats_m^{(i)}}}{2\sigma^2}\right).
  \label{eq:mle}
\end{equation}
The solution yields soft regime responsibilities
$\gamma_{im}^* = p(z_i{=}m \mid s^{(i)},\Phi^*,\pi^*,\sigma^{2*})$
that serve as supervision for Phase~2.



\paragraph{Phase 2: Online operator learning.}

With bases $\{\Phi_m^*\}$ and responsibilities $\{\gamma_{im}^*\}$ frozen,
we train a gating network (\emph{GatingNet}) and $M$ branch networks (\emph{BranchNets}) that map a new input
$(u_0, \kappa)$ to a weighted combination of local regime predictions.
The precise architecture and loss are given in Section~\ref{sec:online}.
The quality of the final operator is measured by the mean relative $L^2$ error:
\begin{equation}
  \varepsilon \;=\; \frac{1}{N_\mathrm{test}}
  \sum_{i=1}^{N_\mathrm{test}}
  \frac{\norm{s^{(i)} - \hat{s}(\cdot;\,u_0^{(i)},\kappa^{(i)})}}
       {\norm{s^{(i)}}} .
  \label{eq:metric}
\end{equation}















\section{Method}\label{sec:theory}



All methods studied in this work share a common two-phase structure:
an \emph{offline} basis-fitting phase that constructs a low-rank
approximation to the solution manifold from snapshots in $\mathcal{D}$,
and an \emph{online} operator-learning phase that trains a branch
network to predict expansion coefficients from a new input
$(u_0, \kappa)$.
The four methods differ along two independent axes:

\begin{itemize}
  \item \textbf{Basis type.}
    A \emph{linear} (POD) basis uses SVD of the (weighted) snapshot
    matrix; a \emph{nonlinear} (NeuralPOD) basis learns spatiotemporal
    mode networks by sequential residual minimization.
  \item \textbf{Number of regimes.}
    A \emph{single-regime} operator fits one global basis to all
    snapshots; a \emph{multi-regime} operator (TEMPO) discovers $M$
    latent regimes via EM and assigns a dedicated basis to each.
\end{itemize}

\noindent
This gives four combinations:
POD-DeepONet~\cite{Lu2022} (linear, global),
NeuralPOD-DeepONet (nonlinear, global, our extension of~\cite{Mou2026}),
TEMPO(POD) (linear, $M$ regimes), and
TEMPO(NeuralPOD) (nonlinear, $M$ regimes).


\subsection{Basis Representations}\label{sec:bases}


\subsubsection{POD Basis}

Let $\phi_0 := \bar{s} = \tfrac{1}{N}\sum_i s^{(i)}$ be the mean field
and $\tilde{S} \in \mathbb{R}^{N_y \times N}$ the matrix with columns
$\tilde{s}^{(i)} = s^{(i)} - \phi_0$.
The rank-$K$ truncated SVD $\tilde{S} \approx U_K \Sigma_K V_K^\top$
yields orthonormal modes $\{\phi_k\}_{k=1}^K$ (columns of $U_K$) and
coefficients $\lambda_k^{(i)} = [\Sigma_K V_K^\top]_{ki}$; the
approximation
\begin{equation}
  \hat{s}^{(i)} = \phi_0 + \sum_{k=1}^{K} \lambda_k^{(i)}\,\phi_k
  \label{eq:pod_approx}
\end{equation}
is optimal in the $L^2$ sense among all rank-$K$ affine representations.

For TEMPO(POD), the M-step~\eqref{eq:mstep_phi} is solved analytically:
set $\phi_0^{(m)} := \bar{s}_m = \sum_i \gamma_{im} s^{(i)}$ and apply SVD to
the weighted centered matrix with columns
$\sqrt{\gamma_{im}}\,(s^{(i)}-\phi_0^{(m)})$.
This is the closed-form linear analogue of Algorithm~\ref{alg:npod}.

\subsubsection{NeuralPOD Basis}\label{sec:neuralpod_basis}

To overcome the linearity constraint, we parametrize each mode
function $\phi_k : \Omega{\times}\mathcal{T} \to \mathbb{R}$ as a
Fourier feature network~\cite{Tancik2020}; we call the resulting basis a \emph{Fourier NeuralPOD} basis:
\begin{equation}
  \phi_k(y) \;=\; f_{\theta_k}\!\bigl(\psi(y)\bigr),
  \qquad
  \psi(y) \;=\;
  \bigl[\sin(2\pi By),\;\cos(2\pi By)\bigr],
  \label{eq:ffn}
\end{equation}
where $B \in \mathbb{R}^{F \times d}$ is a fixed random matrix drawn
once at initialization and $f_{\theta_k} : \mathbb{R}^{2F} \to \mathbb{R}$
is a small MLP with Tanh activations.
Plain MLPs exhibit spectral bias~\cite{Rahaman2019}; the Fourier encoding
resolves this.
For multi-scale coverage, the rows of $B$ are partitioned evenly across
$L$ frequency scales $\sigma_1 < \cdots < \sigma_L$: the $\ell$-th block is drawn from
$\mathcal{N}(0,\,\sigma_\ell^2 I)$.
The mean field $\phi_0$ uses the same architecture.

The scalar amplitudes $\lambda_k^{(i)} \in \mathbb{R}$ are direct
learnable parameters, one per training trajectory, optimized jointly
with $\phi_k$.
After Phase~1 they are frozen and serve as regression targets for the
BranchNet in Phase~2.

Modes are learned by \emph{greedy residual minimization}~\cite{Mou2026}.
Let $r^{(0,i)} = s^{(i)} - \phi_0$, where $\phi_0$ is fitted to the snapshot mean $\bar{s}$ defined above.
For each subsequent mode $k \ge 1$, given the residual
\begin{equation}
  r^{(k-1,i)}(y)
  = s^{(i)}(y)
  - \phi_0(y)
  - \sum_{\ell=1}^{k-1}\lambda_\ell^{(i)}\,\phi_\ell(y),
  \label{eq:residual_global}
\end{equation}
we solve
\begin{equation}
  \min_{\phi_k,\,\{\lambda_k^{(i)}\}}
  \sum_{i=1}^{N} \sum_j w_j
  \Bigl(\lambda_k^{(i)}\,\phi_k(y_j) - r_j^{(k-1,i)}\Bigr)^{\!2}
  \label{eq:neuralpod_loss}
\end{equation}
(the global case with uniform weights; the $\gamma$-weighted
generalisation used in TEMPO is Algorithm~\ref{alg:npod}).
Mode addition stops when the residual falls below fraction $\tau$ of
the initial variance,
\begin{equation}
  \sum_{i=1}^N \|r^{(k,i)}\|_w^2
  \;<\;
  \tau \cdot \sum_{i=1}^N \|s^{(i)} - \bar{s}\|_w^2,
  \label{eq:npod_tol}
\end{equation}
or $k = K_{\max}$.
Each mode captures the largest remaining variance component — the nonlinear analogue of Gram--Schmidt orthogonalization.

\subsection{Single-Regime Operators}\label{sec:single}


\subsubsection{POD-DeepONet}

POD-DeepONet~\cite{Lu2022} replaces the learned trunk network of
DeepONet~\cite{Lu2019} with the fixed POD basis.
After computing the global modes $\{\phi_k\}_{k=1}^K$ and training
coefficients $\{\lambda^{(i)}\}$ offline, a branch network
$\mathcal{B}_\theta : \mathbb{R}^{m_s + d_\kappa} \to \mathbb{R}^K$
is trained online by regression against these coefficients:
\begin{equation}
  \mathcal{L}_{\mathrm{branch}}
  = \frac{1}{N}\sum_{i=1}^{N}
    \bigl\|\mathcal{B}_\theta(u_0^{(i)},\,\kappa^{(i)})
    - \lambda^{(i)}\bigr\|^2.
  \label{eq:pod_branch_loss}
\end{equation}
The prediction at any $(x,t) \in \Omega \times \mathcal{T}$ is
\begin{equation}
  \hat{s}(x,t;\,u_0,\kappa)
  = \phi_0(x,t)
  + \sum_{k=1}^{K}
    \mathcal{B}_{\theta,k}(u_0,\kappa)\,\phi_k(x,t).
  \label{eq:pod_predict}
\end{equation}

\subsubsection{NeuralPOD-DeepONet}\label{sec:neuralpod_deeponet}

Mou et al.~\cite{Mou2026} introduce NeuralPOD as a dimensionality
reduction for PDE solutions, establishing the greedy basis-fitting
procedure of Section~\ref{sec:neuralpod_basis}.
However, they do not construct a deployable solution operator: no
network is proposed to predict the mode coefficients for an unseen
initial condition.
We close this gap by substituting the NeuralPOD basis into the
POD-DeepONet framework; we call this combination \emph{NeuralPOD-DeepONet}.

The procedure follows POD-DeepONet exactly with the Fourier NeuralPOD basis substituted for POD.
Offline fitting yields coefficients $\boldsymbol{\lambda}^{(i)} \in \mathbb{R}^K$; online, a branch network minimizes~\eqref{eq:pod_branch_loss} and prediction follows~\eqref{eq:pod_predict}.
At equal mode count $K$, the nonlinear basis spans a strictly richer function space than POD.


\subsection{TEMPO: Offline Regime Discovery}\label{sec:offline}


A single global basis — whether linear or nonlinear — must compromise
between qualitatively distinct regimes, as the optimal basis geometry
differs fundamentally between them.
TEMPO avoids this by constructing a dedicated basis for each of $M$
latent regimes, discovered from unlabelled trajectories via EM on
the generative model of Definition~\ref{def:generative}.

\paragraph{Initialization.}
A global POD is computed on all $N$ trajectories; each trajectory is
projected to a $P$-dimensional coefficient vector
$\alpha^{(i)} \in \mathbb{R}^P$.
Running a standard GMM on $\{\alpha^{(i)}\}$ yields initial
responsibilities $\gamma_{im}^{(0)}$, weights $\pi_m^{(0)}$,
and per-regime statistics $(\mu_m^{(0)}, \Sigma_m^{(0)})$.
Each regime basis $\Phi_m^{(0)}$ is then fitted to the
responsibility-weighted ensemble: for TEMPO(NeuralPOD) via
\textsc{WeightedNeuralPOD} (Algorithm~\ref{alg:npod}); for TEMPO(POD)
via the weighted SVD of Section~\ref{sec:bases}.
The noise level $\sigma^2$ is calibrated from the initial
reconstructions rather than treated as a free hyperparameter:
\begin{equation}
  \sigma^2 \;\leftarrow\;
  \frac{1}{2N}\sum_{i=1}^N\sum_{m=1}^M
  \gamma_{im}^{(0)}\,\bigl\|s^{(i)} - \hat{s}_m^{(i)}\bigr\|_w^2,
  \label{eq:sigma2_calib}
\end{equation}
which ties the softness of the E-step directly to the quality of the
initial bases.

\paragraph{E-step.}
At iteration $q$, the responsibility of trajectory $i$ for regime $m$
is updated by Bayes' rule on the Gaussian likelihood~\eqref{eq:likelihood}:
\begin{equation}
  \gamma_{im}^{(q)}
  = \frac{
    \pi_m^{(q-1)}\,
    \exp\!\Bigl(-\dfrac{\|s^{(i)}-\hat{s}_m^{(i)}\|_w^2}{2\sigma^2}\Bigr)
  }{
    \sum_{j=1}^{M}
    \pi_j^{(q-1)}\,
    \exp\!\Bigl(-\dfrac{\|s^{(i)}-\hat{s}_j^{(i)}\|_w^2}{2\sigma^2}\Bigr)
  },
  \label{eq:estep}
\end{equation}
where $\hat{s}_m^{(i)}$ is evaluated via~\eqref{eq:separable}
using $\Phi_m^{(q-1)}$ and $\Lambda_m^{(q-1)}$.
As the bases improve in the M-step, the reconstructions $\hat{s}_m^{(i)}$
sharpen and the E-step automatically refines regime assignments in response —
distinguishing TEMPO from a GMM applied to a fixed feature space.
When $\kappa$ drives large amplitude variation (as with $\beta\in[0.01,100]$
in the Darcy experiments), replacing the squared norm with the relative squared norm
$\|s^{(i)}-\hat{s}_m^{(i)}\|_w^2/\|s^{(i)}\|_w^2$ in both~\eqref{eq:estep}
and~\eqref{eq:sigma2_calib} gives a scale-invariant analogue.
When $\kappa$ is observed, the GMM initialization can operate in $\log\kappa$
space, and an optional $\kappa$-prior term in the E-step anchors assignments
to the known parameter structure.

\paragraph{M-step.}
With responsibilities fixed, the M-step gives a closed-form update
for the mixture weights,
\begin{equation}
  \pi_m^{(q)} \leftarrow \frac{N_m^{(q)}}{N},
  \qquad N_m^{(q)} = \sum_{i=1}^{N}\gamma_{im}^{(q)},
  \label{eq:mstep_pi}
\end{equation}
and a weighted reconstruction problem for each basis:
\begin{equation}
  \min_{\Phi_m,\,\Lambda_m}
  \sum_{i=1}^{N} \gamma_{im}^{(q)}\,\|s^{(i)} - \hat{s}_m^{(i)}\|_w^2,
  \label{eq:mstep_phi}
\end{equation}
solved by weighted SVD for TEMPO(POD) or
\textsc{WeightedNeuralPOD} (Algorithm~\ref{alg:npod}) for
TEMPO(NeuralPOD).
Since gradient descent solves the NeuralPOD subproblem without guaranteeing
global optimality, TEMPO is a Generalized EM algorithm~\cite{Neal1998},
which guarantees that the observed-data log-likelihood is non-decreasing
at every iteration.

\paragraph{Adaptive basis update.}
Retraining all $M$ NeuralPOD bases at every EM iteration is
computationally prohibitive.
We introduce a distribution-shift criterion to determine whether each
basis requires updating.
Let $\alpha^{(i)}$ denote the global POD projection of $s^{(i)}$,
computed once during initialization and fixed thereafter.
The responsibility-weighted first and second moments of regime $m$ are
\begin{align}
  \mu_m^{(q)}
  &= \frac{1}{N_m^{(q)}}\sum_{i=1}^N
    \gamma_{im}^{(q)}\,\alpha^{(i)},
  \label{eq:mu}\\
  \Sigma_m^{(q)}
  &= \frac{1}{N_m^{(q)}}\sum_{i=1}^N
    \gamma_{im}^{(q)}\,
    (\alpha^{(i)}-\mu_m^{(q)})(\alpha^{(i)}-\mu_m^{(q)})^\top.
  \label{eq:Sigma}
\end{align}
The distribution-shift criterion measures the relative change in these
moments between consecutive iterations:
\begin{equation}
  \Delta_m^{(q)}
  = \frac{\|\mu_m^{(q)} - \mu_m^{(q-1)}\|_2}
         {\|\mu_m^{(q-1)}\|_2 + \epsilon}
  + \frac{\|\Sigma_m^{(q)} - \Sigma_m^{(q-1)}\|_F}
         {\|\Sigma_m^{(q-1)}\|_F + \epsilon}.
  \label{eq:Delta}
\end{equation}
\noindent where $\epsilon > 0$ is a small numerical stabiliser.
$\Delta_m^{(q)}$ measures how much the composition of regime $m$ has changed,
not just the volume of reassignments.
Two trajectories with very different $\alpha^{(i)}$ swapping regimes
produce a large $\Delta_m^{(q)}$ and trigger a basis update;
two similar snapshots doing so do not.
Given thresholds $0 < \varepsilon_s < \varepsilon_l$, the update
decision at iteration $q$ follows a three-way rule:

\begin{center}
\renewcommand{\arraystretch}{1.3}\small
\begin{tabular}{lll}
\toprule
Condition & Action & Rationale \\
\midrule
$\Delta_m^{(q)} < \varepsilon_s$
  & \textsc{Skip}
  & Data manifold unchanged; basis remains accurate \\
$\varepsilon_s \le \Delta_m^{(q)} < \varepsilon_l$
  & \textsc{Incremental}
  & Moderate drift; add or prune one mode \\
$\Delta_m^{(q)} \ge \varepsilon_l$
  & \textsc{Full rerun}
  & Large shift; rebuild from random initialization \\
\bottomrule
\end{tabular}
\end{center}

\noindent
In the \textsc{Incremental} case, we compute the residual of the
current basis under the updated responsibilities $\gamma^{(q)}$, add
a new mode if $\sum_i \gamma_{im}^{(q)}\|r_m^{(K_m,i)}\|_w^2 \ge \tau\cdot\sum_i \gamma_{im}^{(q)}\|s^{(i)}-\bar{s}_m^{(q)}\|_w^2$,
and prune the last mode if its weighted contribution falls below
$\varepsilon_p$.
The \textsc{Full rerun} case uses a cold-start random initialization,
as warm-starting may trap gradient descent in the geometry of the previous regime.

\paragraph{Convergence and model selection.}
EM terminates when $\max_m \Delta_m^{(q)} < \varepsilon_c$.
The number of regimes $M \in \{2,3,4\}$ is selected by BIC on the GMM fit,
with mean assignment entropy as a secondary criterion;
see Table~\ref{tab:regime_selection}.
The complete offline procedure is given in Algorithm~\ref{alg:tempo}.

\subsection{TEMPO: Online Gated Operator}\label{sec:online}


After Algorithm~\ref{alg:tempo} converges, all regime bases
$\{\Phi_m^*\}$ and responsibilities $\{\gamma_{im}^*\}$ are frozen.
The online phase trains a GatingNet
$G: \mathbb{R}^{m_s + d_\kappa} \to \Delta^{M-1}$
that maps $(u_0, \kappa)$ to regime weights, and $M$ BranchNets
$\mathcal{B}_m: \mathbb{R}^{m_s + d_\kappa} \to \mathbb{R}^{K_m}$
that predict expansion coefficients for each regime.
The TEMPO prediction at query point $y = (x,t) \in \Omega\times\mathcal{T}$ is
\begin{equation}
  \hat{s}(y;\,u_0,\kappa)
  = \sum_{m=1}^{M} g_m(u_0,\kappa)
    \left[
      \phi_0^{(m)}\!(y)
      + \sum_{k=1}^{K_m}
        b_{m,k}(u_0,\kappa)\;
        \phi_k^{(m)}\!(y)
    \right],
  \label{eq:predict}
\end{equation}
where $\mathbf{g} = G(u_0,\kappa)$ and
$\mathbf{b}_m = \mathcal{B}_m(u_0,\kappa)$.
The BranchNet outputs $\mathbf{b}_m$ predict the expansion coefficients
$\boldsymbol{\lambda}_m^{(i)}$ for new inputs $(u_0,\kappa)$ not seen in training.
The prediction is mesh-free and evaluable at any $y\in\Omega\times\mathcal{T}$
without retraining.

\paragraph{Training objective.}
Writing $g_m^{(i)} = g_m(u_0^{(i)},\kappa^{(i)})$, the online loss is
\begin{equation}
  \mathcal{L}_{\mathrm{online}}
  = \underbrace{
      \frac{1}{N}\sum_{i=1}^{N}
      \|s^{(i)} - \hat{s}(\cdot;\,u_0^{(i)},\kappa^{(i)})\|_w^2
    }_{\mathcal{L}_{\mathrm{data}}}
  +\;\lambda_{\mathrm{KL}}
    \underbrace{
      \frac{1}{N}\sum_{i,m}
      g_m^{(i)}\log\frac{g_m^{(i)}}{\gamma_{im}^*}
    }_{\mathcal{L}_{\mathrm{KL}}}
  -\;\lambda_{H}
    \underbrace{
      \frac{1}{N}\sum_{i,m}
      g_m^{(i)}\log g_m^{(i)}
    }_{-\mathcal{L}_{\mathrm{ent}}}.
  \label{eq:online_loss}
\end{equation}
$\mathcal{L}_{\mathrm{data}}$ drives reconstruction accuracy.
$\mathcal{L}_{\mathrm{KL}}$ anchors the online gating to the
offline regime partition: it penalizes deviation of $\mathbf{g}^{(i)}$
from the EM responsibilities $\boldsymbol{\gamma}_i^*$, so the
network cannot discard the discovered structure in pursuit of a
simpler routing.
$\mathcal{L}_{\mathrm{ent}}$ prevents collapse to a single expert,
ensuring all $M$ regime bases contribute.
The hyperparameters $\lambda_{\mathrm{KL}}, \lambda_H \ge 0$ trade
off these objectives.







\section{Computational Experiments}
\label{sec:exp}

We evaluate TEMPO on parametric PDE benchmarks, verifying three claims:
(i)~the offline EM phase recovers interpretable, physics-consistent regime
structure from unlabelled snapshots;
(ii)~locally adapted bases achieve lower reconstruction error than a
single global basis at equal mode count;
(iii)~a jointly-trained gated operator with regime-specific bases outperforms
a single jointly-trained basis without regime discovery.

\subsection{Experimental Setup}
\label{sec:setup}

\paragraph{Datasets.}
We use two parametric PDE benchmarks from PDEBench~\cite{Takamoto2022}.

\textbf{1D Burgers' equation.}
\begin{equation}
  u_t + u\,u_x = \nu\,u_{xx},
  \qquad x\in(-1,1),\quad t\in(0,2],
  \label{eq:burgers}
\end{equation}
with periodic boundary conditions and random initial conditions.
Viscosity $\nu\in\{0.001,\,0.01,\,0.1,\,1.0\}$ spans four qualitatively
distinct solution regimes: smooth diffusive profiles at $\nu{=}1.0$ and
near-discontinuous shock structures at $\nu{=}0.001$.
Each value provides $N{=}9{,}500$ trajectories on $N_x{=}1{,}024$
spatial points and $N_t{=}201$ time steps ($38{,}000$ total).

\textbf{2D Darcy flow.}
\begin{equation}
  -\nabla\cdot\bigl(a(x)\,\nabla u(x)\bigr) = \beta,
  \qquad x\in(0,1)^2,
  \label{eq:darcy}
\end{equation}
with zero Dirichlet boundary conditions.
Here $a(x)$ is a spatially varying permeability field sampled from a
random sum of Gaussians, and $\beta\in\{0.01,\,0.1,\,1.0,\,10.0,\,100.0\}$
is the constant source term.
Each value provides $N{=}8{,}000$ training and $1{,}000$ test samples
discretised on a $128{\times}128$ uniform grid ($N_y{=}16{,}384$ spatial points),
$40{,}000$ trajectories in total.
The operator maps $a(x)$ to the pressure solution $u(x)$.

\paragraph{Methods.}
All methods are implemented and trained by us.
POD-DeepONet~\cite{Lu2022} and NeuralPOD-DeepONet (\textit{ours}) are trained
one model per $\nu$ (or $\beta$) value and evaluated in-distribution;
they serve as a reference ceiling for jointly-trained approaches.
POD-DeepONet trained jointly on all parameter values without regime discovery
is the primary baseline for TEMPO.
TEMPO(POD) uses responsibility-weighted SVD per regime;
TEMPO(NeuralPOD) uses Fourier NeuralPOD bases.
Both TEMPO variants train a single model on all parameter values jointly,
taking $(u_0, \kappa)$ as input at inference.

\paragraph{Metric.}
All methods are evaluated by the mean relative $L^2$ error $\varepsilon$
defined in~\eqref{eq:metric}, computed per parameter value and reported
in Tables~\ref{tab:results}--\ref{tab:darcy}.

%==========================================================================
\subsection{Main Results}
\label{sec:results}
%==========================================================================

\paragraph{1D Burgers' equation.}
Table~\ref{tab:results} reports the full cross-viscosity error matrix for the
Burgers' benchmark.
Each per-$\nu$ specialist achieves low error in-distribution (diagonal entries),
but degrades sharply outside its training regime: for example, a model trained
on $\nu{=}0.1$ reaches $0.048$ in-distribution yet $0.461$ on $\nu{=}0.001$,
a ten-fold increase.
This cross-regime failure motivates joint training with automatic regime
discovery.
TEMPO is compared against a jointly-trained POD-DeepONet without regime
discovery; TEMPO improves on this baseline by learning regime-specific bases.

\paragraph{2D Darcy flow.}
Table~\ref{tab:darcy} shows the cross-$\beta$ error matrix for Darcy flow.
Per-$\beta$ specialists fail even more severely than in Burgers: the Darcy
solution scales linearly with~$\beta$, so a model trained on one value
predicts completely wrong amplitudes on others.
Errors grow by roughly one order of magnitude per decade in~$\beta$,
reaching ${>}1000$ for the most distant pair.
TEMPO trains jointly on $\beta\in\{0.1,1,10,100\}$ and recovers reasonable
accuracy across all regimes.

% TABLE 1: Burgers
\begin{table}[t]
\centering
\small
\setlength{\tabcolsep}{5pt}
\caption{\textbf{Relative $L^2$ test error on 1D Burgers' equation.}
  Each per-$\nu$ specialist is trained on a single viscosity and evaluated on
  all four; \underline{underlined} entries are in-distribution.
  The large off-diagonal errors motivate joint training with regime discovery.
  TEMPO trains jointly on $\nu\in\{0.001,0.1,1.0\}$; the $\nu{=}0.01$ column
  is omitted for jointly-trained methods~(\textemdash).
  \textbf{Bold}: best per column among jointly-trained methods.}
\label{tab:results}
\begin{tabular}{l c c c c c}
\toprule
Method & Train~$\nu$
  & $\nu{=}0.001$ & $\nu{=}0.01$ & $\nu{=}0.1$ & $\nu{=}1.0$ \\
\midrule
\multicolumn{6}{l}{\textit{Per-$\nu$ specialists (single-viscosity training)}} \\[2pt]
\multirow{4}{*}{POD-DeepONet~\cite{Lu2022}}
  & 0.001 & \underline{0.273} & 0.215 & 0.297 & 0.639 \\
  & 0.01  & 0.274 & \underline{0.214} & 0.274 & 0.614 \\
  & 0.1   & 0.349 & 0.299 & \underline{0.048} & 0.344 \\
  & 1.0   & 0.461 & 0.425 & 0.260 & \underline{0.025} \\[4pt]
\multirow{4}{*}{NeuralPOD-DeepONet~\textit{(ours)}}
  & 0.001 & \underline{0.392} & 0.352 & 0.288 & 0.535 \\
  & 0.01  & 0.391 & \underline{0.351} & 0.284 & 0.541 \\
  & 0.1   & 0.426 & 0.388 & \underline{0.229} & 0.371 \\
  & 1.0   & 0.483 & 0.451 & 0.315 & \underline{0.183} \\
\midrule
\multicolumn{6}{l}{\textit{Joint training ($\nu\in\{0.001,0.1,1.0\}$)}} \\[2pt]
POD-DeepONet~\cite{Lu2022}
  & all & \textemdash & \textemdash & \textemdash & \textemdash \\
TEMPO(POD)~\textit{(ours)}
  & $M{=}3$ & \textbf{0.317} & \textemdash & \textbf{0.168} & \textbf{0.182} \\
TEMPO(NeuralPOD)~\textit{(ours)}
  & $M{=}3$ & 0.571 & \textemdash & 0.502 & 0.376 \\
\bottomrule
\end{tabular}
\end{table}

% TABLE 2: Darcy
\begin{table}[t]
\centering
\small
\setlength{\tabcolsep}{4pt}
\caption{\textbf{Relative $L^2$ test error on 2D Darcy flow.}
  Each per-$\beta$ specialist is trained on a single source value and evaluated
  on all five; \underline{underlined} entries are in-distribution.
  Off-diagonal errors exceed 1 because the Darcy solution scales with~$\beta$,
  so a specialist trained on~$\beta{=}100$ predicts amplitudes ${\sim}10^{4}$
  times larger than the ground truth at~$\beta{=}0.01$.
  TEMPO trains jointly on $\beta\in\{0.1,1,10,100\}$; the $\beta{=}0.01$ regime
  is excluded from joint training~(\textemdash).
  \textbf{Bold}: best per column among jointly-trained methods.}
\label{tab:darcy}
\begin{tabular}{l c c c c c c}
\toprule
Method & Train~$\beta$
  & $\beta{=}0.01$ & $\beta{=}0.1$ & $\beta{=}1.0$
  & $\beta{=}10$ & $\beta{=}100$ \\
\midrule
\multicolumn{7}{l}{\textit{Per-$\beta$ specialists (single-value training)}} \\[2pt]
\multirow{5}{*}{POD-DeepONet~\cite{Lu2022}}
  & 0.01  & \underline{0.449} & 0.791 & 0.961 & 0.996 & ${>}1$ \\
  & 0.1   & ${>}1$   & \underline{0.148} & 0.871 & 0.986 & 0.999 \\
  & 1.0   & ${>}10$  & ${>}1$   & \underline{0.048} & 0.897 & 0.990 \\
  & 10    & ${>}100$ & ${>}10$  & ${>}1$   & \underline{0.035} & 0.900 \\
  & 100   & ${>}1000$& ${>}100$ & ${>}10$  & ${>}1$   & \underline{0.042} \\
\midrule
\multicolumn{7}{l}{\textit{Joint training ($\beta\in\{0.1,1,10,100\}$)}} \\[2pt]
POD-DeepONet~\cite{Lu2022}
  & all
  & \textemdash & \textemdash & \textemdash & \textemdash & \textemdash \\
TEMPO(POD)~\textit{(ours)}
  & $M{=}5$
  & \textemdash & \textbf{0.775} & \textbf{0.178} & \textbf{0.136} & \textbf{0.073} \\
TEMPO(NeuralPOD)~\textit{(ours)}
  & $M{=}5$
  & \textemdash & 3.813 & 0.549 & 0.322 & 0.311 \\
\bottomrule
\end{tabular}
\end{table}

%==========================================================================
\subsection{Regime Discovery}
\label{sec:regime_exp}
%==========================================================================

\paragraph{Number of regimes.}
We sweep $M\in\{2,3,4\}$ and report three criteria in
Table~\ref{tab:regime_selection}: the EM log-likelihood (BIC, higher is better),
mean assignment entropy $H = -\frac{1}{N}\sum_{i,m}\gamma_{im}\log\gamma_{im}$,
and overall test error~$\bar{\varepsilon}$.
We select $M{=}3$ based on the entropy elbow: the gain from $M{=}2$ to $M{=}3$
is $+0.30$, while the gain from $M{=}3$ to $M{=}4$ drops to $+0.18$,
indicating that the third regime captures a genuine structure that a fourth
does not.
The UMAP embedding (Figure~\ref{fig:umap_em}) confirms three distinct manifold
branches, one per discovered regime.

% TABLE 3: Regime selection
\begin{table}[h]
\centering
\small
\setlength{\tabcolsep}{12pt}
\caption{Regime count selection on 1D Burgers'.
  BIC: EM log-likelihood (higher is better);
  $H$: mean assignment entropy;
  $\bar{\varepsilon}$: mean relative $L^2$ test error.}
\label{tab:regime_selection}
\begin{tabular}{c c c c}
\toprule
$M$ & BIC (${\times}10^{4}$) & $H$ & $\bar{\varepsilon}$ \\
\midrule
2          & 3.90 & 0.622 & 0.209 \\
\textbf{3} & \textbf{4.92} & \textbf{0.920} & \textbf{0.222} \\
4          & 5.47 & 1.099 & 0.228 \\
\bottomrule
\end{tabular}
\end{table}

\paragraph{Gating structure.}
Table~\ref{tab:gating} reports the mean EM assignment weight
$\bar{\gamma}_m^\nu = |\mathcal{T}_\nu|^{-1}\sum_{i\in\mathcal{T}_\nu}\gamma_{im}^*$
for each viscosity class, alongside the global mixture weights~$\pi_m$.
The per-$\nu$ weights are nearly identical to~$\pi_m$ across all three regimes,
meaning each viscosity class is split between regimes in almost exactly the
same proportion as the full dataset.
This confirms that regime membership is determined by the shape of the initial
condition~$u_0$, not by the viscosity label~$\nu$ directly, consistent with
the UMAP structure in Figure~\ref{fig:umap_em}.

% TABLE 4: Gating weights
\begin{table}[h]
\centering
\small
\setlength{\tabcolsep}{14pt}
\caption{Mean EM assignment weights per viscosity (1D Burgers',
  TEMPO(POD), $M{=}3$) compared to global mixture weights~$\pi_m$.
  Each data row sums to~$1$.
  The per-$\nu$ weights closely match~$\pi_m$, confirming that regime
  membership is driven by the shape of~$u_0$, not by~$\nu$.}
\label{tab:gating}
\begin{tabular}{c c c c}
\toprule
$\nu$ & Regime~1 & Regime~2 & Regime~3 \\
\midrule
0.001 & 0.32 & 0.28 & 0.40 \\
0.1   & 0.32 & 0.26 & 0.42 \\
1.0   & 0.33 & 0.26 & 0.41 \\
\midrule
$\pi_m$ & 0.32 & 0.27 & 0.41 \\
\bottomrule
\end{tabular}
\end{table}

\paragraph{Regime structure.}
Figure~\ref{fig:umap_em} shows the UMAP embedding of all training snapshots,
coloured by regime assignment and by~$\nu$.
The three discovered regimes do \emph{not} coincide with individual viscosity
values: each manifold branch spans multiple~$\nu$, confirming that the EM
partition reflects intrinsic solution geometry rather than the physical parameter label.
After convergence, mixture weights are
$\boldsymbol{\pi}{=}(0.324,\,0.267,\,0.409)$, yielding a balanced,
non-degenerate partition.

\begin{figure}[t]
\centering
\includegraphics[width=\textwidth]{TEMPO/main/pictures/TEMPO_phase1.png}
\caption{TEMPO regime structure on 1D Burgers' ($M{=}3$, 19 EM iterations).
  \textit{Left:} UMAP coloured by regime assignment.
  \textit{Centre:} same embedding coloured by~$\nu$; each branch spans
  multiple viscosity values.
  \textit{Right:} GMM fit in POD coefficient space $(\alpha_1,\alpha_2)$;
  stars mark regime centroids~$\mu_m$.}
\label{fig:umap_em}
\end{figure}

%==========================================================================
\subsection{Reconstruction Quality}
\label{sec:recon_exp}
%==========================================================================

Figure~\ref{fig:reconstruct} shows a representative space--time prediction
from TEMPO(NeuralPOD) on 1D Burgers' ($\nu{=}1.0$).
The prediction closely tracks the ground truth with small, uniformly distributed
residuals and no systematic error near the diffusive front.
Additional reconstructions and error distributions for all baselines are
provided in Appendix~\ref{app:exp}.

\begin{figure}[h]
\centering
\includegraphics[width=\textwidth]{TEMPO/main/pictures/reconstruct.png}
\caption{Space--time reconstruction on 1D Burgers' ($\nu{=}1.0$,
  TEMPO(NeuralPOD)): ground truth (\textit{left}), prediction
  (\textit{centre}), and point-wise absolute error (\textit{right}).}
\label{fig:reconstruct}
\end{figure}



\section{Conclusion}

Operator learning typically uses a single global representation across the full parameter space.
When the governing parameter drives qualitative change in solution structure, a single basis must fit incompatible geometries simultaneously, concentrating error where accurate prediction matters most.

TEMPO addresses this by treating regime membership as a latent variable.
EM discovers $M$ regimes from unlabelled snapshots and fits a dedicated basis to each, without regime supervision.
On 1D Burgers' across three orders of viscosity, the discovered regimes cross viscosity boundaries, confirming that the partition reflects solution geometry rather than the physical parameter directly.
Regime-specific bases reduce reconstruction error relative to a global basis at equal mode count.
The gated online operator outperforms a single joint baseline across the full parameter range.

Several directions remain open.
Conditioning the basis functions on $\kappa$ would allow each mode to encode parameter-dependent structure directly.
Extending TEMPO to higher-dimensional benchmarks will test whether regime discovery generalises beyond one-dimensional shock dynamics.
Convergence guarantees for Generalized EM under nonlinear basis parameterisation remain an open theoretical problem.

\bibliographystyle{unsrtnat}
\bibliography{references}


\newpage
\appendix
\section{Appendix}\label{app}

\subsection{Algorithms}\label{app:alg}

\begin{algorithm}[H]
\caption{\textsc{WeightedNeuralPOD}$\bigl(
  \{s^{(i)}, \gamma_{im}\}_{i=1}^N,\;
  \tau\bigr)$}
\label{alg:npod}
\begin{algorithmic}[1]
\State $\bar{s}_m \leftarrow \displaystyle\sum_{i=1}^N \gamma_{im}\,s^{(i)}$
  \hfill\Comment{responsibility-weighted mean trajectory}
\State Train $\phi_0^{(m)}$ by minimizing
  $\displaystyle\sum_j w_j\bigl(\phi_0^{(m)}(y_j) - (\bar{s}_m)_j\bigr)^2$
\State $r^{(0,i)} \leftarrow s^{(i)} - \phi_0^{(m)}(y)$
  \enspace for all $i$;\enspace $k \leftarrow 0$
\State $V_m \leftarrow \displaystyle\sum_{i=1}^N \gamma_{im}\,\|s^{(i)}-\bar{s}_m\|_w^2$
  \hfill\Comment{initial unmodeled variance}
\While{$\displaystyle\sum_{i=1}^N \gamma_{im}\,\|r^{(k,i)}\|_w^2
  \;\ge\; \tau\, V_m$ \textbf{ and } $k < K_{\max}$}
  \State Solve $\gamma_{im}$-weighted~\eqref{eq:neuralpod_loss} for $\phi_{k+1}^{(m)}$ and $\bigl\{\lambda_{k+1}^{(m,i)}\bigr\}$
  \State $r^{(k+1,i)} \leftarrow r^{(k,i)}
    - \lambda_{k+1}^{(m,i)}\,\phi_{k+1}^{(m)}(y)$
    \enspace for all $i$
  \State $k \leftarrow k + 1$
\EndWhile
\State \Return $K_m \!\leftarrow\! k$,\enspace
  $\phi_0^{(m)}$,\enspace
  $\bigl\{\phi_k^{(m)},\,\{\lambda_k^{(m,i)}\}_{i=1}^N\bigr\}_{k=1}^{K_m}$
\end{algorithmic}
\end{algorithm}

\begin{algorithm}[H]
\caption{TEMPO --- Offline Phase}
\label{alg:tempo}
\begin{algorithmic}[1]
\Require $\mathcal{D} = \{u_0^{(i)}, \kappa^{(i)}, s^{(i)}\}_{i=1}^N$;\;
  $M,\;\varepsilon_s,\;\varepsilon_l,\;\varepsilon_p,\;\varepsilon_c,\;\tau$
\Ensure Regime bases $\{\Phi_m^*\}$,\;
  amplitudes $\{\Lambda_m^*\}$,\;
  responsibilities $\{\gamma_{im}^*\}$,\;
  weights $\{\pi_m^*\}$
%
\Statex \textit{Initialization}
\State Compute global POD on $\{s^{(i)}\}$;
  project to $\alpha^{(i)} \in \mathbb{R}^P$ for all $i$
\State Run GMM-EM on $\{\alpha^{(i)}\}$ $\to$
  $\gamma_{im}^{(0)},\;\pi_m^{(0)},\;\mu_m^{(0)},\;\Sigma_m^{(0)}$
\For{$m = 1,\ldots,M$}
  \State $\Phi_m^{(0)},\,\Lambda_m^{(0)} \leftarrow$
    \Call{WeightedNeuralPOD}
    {$\{s^{(i)},\gamma_{im}^{(0)}\},\,\tau$}
    \hfill\Comment{or weighted SVD for TEMPO(POD), see \S\ref{sec:bases}}
\EndFor
\State $\sigma^2 \leftarrow (2N)^{-1}\!\sum_{i,m}
  \gamma_{im}^{(0)}\,\|s^{(i)}{-}\hat{s}_m^{(i)}\|_w^2$
  \hfill\eqref{eq:sigma2_calib}
%
\Statex \textit{EM iterations}
\Repeat
  \Statex \quad\textit{E-step}
  \For{all $i = 1,\ldots,N$ and $m = 1,\ldots,M$}
    \State Compute $\hat{s}_m^{(i)}$ via~\eqref{eq:separable}
      using $\Phi_m^{(q-1)}$ and $\Lambda_m^{(q-1)}$
    \State $\gamma_{im}^{(q)} \leftarrow$ Eq.~\eqref{eq:estep}
  \EndFor
  \Statex \quad\textit{M-step: mixture weights}
  \State $\pi_m^{(q)} \leftarrow N_m^{(q)}/N$
    for all $m$ \hfill\eqref{eq:mstep_pi}
  \Statex \quad\textit{M-step: bases}
  \For{$m = 1,\ldots,M$}
    \State Compute $\mu_m^{(q)}$~\eqref{eq:mu},\;
      $\Sigma_m^{(q)}$~\eqref{eq:Sigma},\;
      $\Delta_m^{(q)}$~\eqref{eq:Delta}
    \If{$\Delta_m^{(q)} < \varepsilon_s$}
      \State $\Phi_m^{(q)} \leftarrow \Phi_m^{(q-1)}$ \Comment{Skip}
    \ElsIf{$\Delta_m^{(q)} < \varepsilon_l$}
      \State Recompute $r_m^{(K_m,i)}$ and $V_m^{(q)} \leftarrow \textstyle\sum_i\gamma_{im}^{(q)}\|s^{(i)}{-}\bar{s}_m^{(q)}\|_w^2$ \Comment{Incremental}
      \If{$\sum_i\gamma_{im}^{(q)}\|r_m^{(K_m,i)}\|_w^2 \ge \tau\cdot V_m^{(q)}$}
        \State Train mode $K_m{+}1$;\; $K_m \leftarrow K_m + 1$
      \EndIf
      \If{$\|\phi_{K_m}^{(m)}\|_w^2\cdot\displaystyle\sum_{i=1}^N
        \gamma_{im}^{(q)}\,\bigl(\lambda_{K_m}^{(m,i)}\bigr)^2 < \varepsilon_p$}
        \State $K_m \leftarrow K_m - 1$
      \EndIf
    \Else
      \State $\Phi_m^{(q)},\,\Lambda_m^{(q)} \leftarrow \Call{WeightedNeuralPOD}{\{s^{(i)},\gamma_{im}^{(q)}\},\,\tau}$ \Comment{Full rerun}
    \EndIf
  \EndFor
\Until{$\max_m \Delta_m^{(q)} < \varepsilon_c$}
\State \Return $\Phi_m^* \!\leftarrow\! \Phi_m^{(q)}$,\;
  $\Lambda_m^* \!\leftarrow\! \Lambda_m^{(q)}$,\;
  $\gamma_{im}^* \!\leftarrow\! \gamma_{im}^{(q)}$,\;
  $\pi_m^* \!\leftarrow\! \pi_m^{(q)}$
\end{algorithmic}
\end{algorithm}

\subsection{Additional Experiments}\label{app:exp}

\paragraph{POD basis quality (1D Burgers').}
Figure~\ref{fig:pod_phase1} shows the singular value spectrum and cumulative
explained variance for the global POD on 1D Burgers'.
$K{=}30$ modes capture $99.99\%$ of the total variance, achieving a
Phase~1 relative $L^2$ reconstruction error of $1.47\%$.

\begin{figure}[h]
\centering
\includegraphics[width=\textwidth]{TEMPO/main/pictures/pod_phase1.png}
\caption{Global POD basis on 1D Burgers' ($\nu{=}1.0$, $K{=}30$ modes).
  \textit{Left:} singular value spectrum.
  \textit{Right:} cumulative explained variance; $K{=}30$ modes attain $99.99\%$.}
\label{fig:pod_phase1}
\end{figure}

\paragraph{NeuralPOD basis quality (1D Burgers').}
Figure~\ref{fig:npod_phase1} shows the greedy residual norm as a function of
mode count; the steep initial decay confirms that the Fourier NeuralPOD basis
concentrates variance efficiently.
Figure~\ref{fig:basis_rec} shows the learned spatial modes alongside
representative reconstructions; modes are smooth and globally supported for
large~$\nu$, and sharply localised near the shock front for small~$\nu$.

\begin{figure}[h]
\centering
\includegraphics[width=\textwidth]{TEMPO/main/pictures/npod_phase1.png}
\caption{Fourier NeuralPOD basis on 1D Burgers' ($\nu{=}1.0$):
  weighted residual norm versus greedy mode count.}
\label{fig:npod_phase1}
\end{figure}

\begin{figure}[h]
\centering
\includegraphics[width=\textwidth]{TEMPO/main/pictures/Neural_POD_rec_basis.png}
\caption{Learned Fourier NeuralPOD basis functions and representative
  reconstructions on 1D Burgers' ($\nu{=}1.0$).}
\label{fig:basis_rec}
\end{figure}

\paragraph{Baseline training curves.}
Figures~\ref{fig:pod_phase2_loss} and~\ref{fig:npod_phase2_loss} show the
Phase~2 MSE training curves for POD-DeepONet and NeuralPOD-DeepONet on
1D Burgers' ($\nu{=}1.0$).

\begin{figure}[h]
\centering
\includegraphics[width=\textwidth]{TEMPO/main/pictures/pod_phase2_loss.png}
\caption{Phase~2 MSE training curve for POD-DeepONet on 1D Burgers' ($\nu{=}1.0$).}
\label{fig:pod_phase2_loss}
\end{figure}

\begin{figure}[h]
\centering
\includegraphics[width=\textwidth]{TEMPO/main/pictures/npod_phase2_loss.png}
\caption{Phase~2 MSE training curve for NeuralPOD-DeepONet on 1D Burgers' ($\nu{=}1.0$).}
\label{fig:npod_phase2_loss}
\end{figure}

\paragraph{Error distributions.}
Figures~\ref{fig:pod_err_dist} and~\ref{fig:npod_err_dist} show the full
test-set error distributions for POD-DeepONet and NeuralPOD-DeepONet on
1D Burgers' ($\nu{=}1.0$).

\begin{figure}[h]
\centering
\includegraphics[width=\textwidth]{TEMPO/main/pictures/pod_err_dist.png}
\caption{POD-DeepONet test-set error distribution on 1D Burgers' ($\nu{=}1.0$).
  \textit{Left:} histogram with mean (red dashed) and median (blue dashed).
  \textit{Right:} sorted relative $L^2$ errors.}
\label{fig:pod_err_dist}
\end{figure}

\begin{figure}[h]
\centering
\includegraphics[width=\textwidth]{TEMPO/main/pictures/npod_err_dist.png}
\caption{NeuralPOD-DeepONet test-set error distribution on 1D Burgers'
  ($\nu{=}1.0$). Layout as in Figure~\ref{fig:pod_err_dist}.}
\label{fig:npod_err_dist}
\end{figure}

\paragraph{Representative reconstructions.}
Figures~\ref{fig:pod_reconstruction} and~\ref{fig:npod_reconstruction} show
representative space--time predictions for POD-DeepONet and
NeuralPOD-DeepONet on 1D Burgers' ($\nu{=}1.0$).

\begin{figure}[h]
\centering
\includegraphics[width=\textwidth]{TEMPO/main/pictures/pod_reconstruction.png}
\caption{POD-DeepONet: representative reconstructions on 1D Burgers'
  ($\nu{=}1.0$). Each row: ground truth (\textit{left}), prediction
  (\textit{centre}), point-wise absolute error (\textit{right}).}
\label{fig:pod_reconstruction}
\end{figure}

\begin{figure}[h]
\centering
\includegraphics[width=\textwidth]{TEMPO/main/pictures/npod_reconstruction.png}
\caption{NeuralPOD-DeepONet: representative reconstructions on 1D Burgers'
  ($\nu{=}1.0$). Layout as in Figure~\ref{fig:pod_reconstruction}.}
\label{fig:npod_reconstruction}
\end{figure}

\end{document}